% Generated by code/make_labor_macros.py from data/examples/invoice_chain.json (version 2026-10-29c); do not edit.
% Numbers of the talent example in part2/theory_labor.tex that come from the running invoice example;
% checked against the JSON by code/verify_part2_labor.py (L3a, L3b, L3f).
\newcommand{\labPremRate}{11\%}
\newcommand{\labTauL}{0.29}
\newcommand{\labWSE}{\$135{,}980}
\newcommand{\labWAI}{\$194{,}649}
\newcommand{\labPremAmt}{\$61{,}600}
\newcommand{\labPremShareF}{2.2\%}
\newcommand{\labLAI}{0.20}
\newcommand{\labLAIp}{2.2\%}
\newcommand{\labContractT}{1.80}
% talent-example numbers from the running invoice example, generated by code/make_labor_macros.py from data/examples/invoice_chain.json
\section{Labor Law, Talent and Jurisdiction}\label{sec:labor}

Sections~\ref{sec:tasks} and \ref{sec:ext} treat labor as a cost per hour and a count of full-time equivalents (FTE). This section adds collective procedures, the response of those who stay, the supply of engineers, and the location of staff and data, and combines them into a bound on pace. Rules are stated as model constraints, as they stood at their primary sources on 29 October 2026; this is not legal advice, and their application depends on facts and case law the model does not represent. Proofs are in Appendix~\ref{app:proofs-labor}, checks in \texttt{code/verify\_part2\_labor.py}.

\subsection{Collective procedures and their lead times}\label{sec:labor-proc}

\begin{assumption}[Jurisdictions and procedures]\label{ass:labor-proc}
Jurisdictions are indexed by $\mathfrak c\in\mathfrak C$ and establishments by $\mathfrak e$, with normal headcount $N^{\rm hc}_{\mathfrak e}$. Time runs from date 0 of Assumption~\ref{ass:tasks-transition}, when the run would start unconstrained.
\begin{enumerate}
\item[(i)] A wave displacing $N^{\rm dis}_{\mathfrak e}$ FTE at $\mathfrak e$ triggers a collective procedure if the dismissals at $\mathfrak e$ within some window of length $T^{\rm win}_{\mathfrak c}$ reach the threshold $N^{\rm thr}_{\mathfrak c}(N^{\rm hc}_{\mathfrak e})$. Covered dismissals cannot take effect before $T^{\rm leg}_{\mathfrak c}:=T^{\rm proc}_{\mathfrak c}+T^{\rm not}_{\mathfrak c}$ after the procedure starts: the statutory minimum for consultation and the authority's wait, plus the part of individual notice that cannot run concurrently. Others need only $T^{\rm not}_{\mathfrak c}$.
\item[(ii)] Each displaced FTE receives $T^{\rm sev}_{\mathfrak c}$ years of its wage (the $T_{\rm sev}$ of Assumption~\ref{ass:tasks-transition} by jurisdiction: statutory floor plus any negotiated plan); a procedure costs $C^{\rm proc}_{\mathfrak c}$.
\item[(iii)] Consultation on the new technology takes $T^{\rm tech}_{\mathfrak c}\ge0$ and can overlap the build of length $T^{\rm bld}$, so the run in $\mathfrak c$ starts at $T^{\rm st}_{\mathfrak c}:=(T^{\rm tech}_{\mathfrak c}-T^{\rm bld})_+$, and a failure is detected, if at all, by the gate $t^{\rm gate}_{\mathfrak c}:=T^{\rm st}_{\mathfrak c}+T_{\rm par}$.
\end{enumerate}
\end{assumption}

\begin{table}[t]
\centering
\caption{Collective procedures as model constraints (verified 29 October 2026)}\label{tab:labor-rules}
\scriptsize
\setlength{\tabcolsep}{3pt}
\begin{tabular}{@{}p{1.65cm}p{4.4cm}p{4.6cm}p{5.15cm}@{}}
\toprule
& Procedure triggered by & Minimum lead before dismissals & Statutory cash floor\\
\midrule
U.S. federal & employer of 100+; at a site, 50+ (closing), or 50+ and 33\%, or 500+ (layoff); 90-day aggregation & 60 days' notice & none; a violation costs up to 60 days' pay and benefits\\
U.S. states & CA: site of 75+, 50+ in 30 days; NY: employer of 50+, 25+ and 33\% or 250+; NJ: employer of 100+, 50+ & 60 days (CA), 90 days (NY, NJ) & NJ: 1 week's pay a year, 4 more if notice is short\\
EU minimum & in 30 days, 10 (21--99), 10\% (100--299), 30 (300+); or 20 in 90 days & consult in good time; 30 days after notifying the authority, extendable to 60 & set nationally\\
Germany & in 30 days, over 5 (21--59), 10\% or over 25 (60--499), 30 (500+) & works-council consultation; hearing before each dismissal; 1 month after notifying the agency (up to 2); notice 1--7 months & none in general; 0.5 month a year if offered; social plan enforceable above thresholds\\
France & firm of 50+ and 10+ in 30 days (job-protection plan) & works council 2, 3 or 4 months (under 100, 100--249, 250+); administration 15 or 21 days; notice 1--2 months & 1/4 month a year to 10 years, 1/3 beyond; reclassification leave up to 12 months at 1,000+\\
U.K. & 20+ at one establishment in 90 days & consultation 30 days (20--99) or 45 days (100+) & 0.5, 1 or 1.5 weeks a year by age, 20 years at most, week capped at \pounds751; protective award up to 180 days' pay\\
India & retrenchment of a worker with a year's service; prior permission only in factories, mines, plantations of 300+ & 1 month's notice or pay; notice to the government & 15 days' pay a completed year, plus 15 days' wages to a re-skilling fund\\
\bottomrule
\end{tabular}

\smallskip
\parbox{\textwidth}{\scriptsize Sources: \citet{us2023warn,cfr2026warn,ca2026warn,nydol2026warn,njdol2026warn,eu1998directive,de2026kschg,de2026betrvg,de2026bgb,fr2026codetravail,uk2026tulrca,uk2026era1996,uk2025era,cipp2026limits,in2020irc,vvgnli2020irc,eu2024regulation,eu2026regulation}. Counts follow each statute's definitions.}
\end{table}

Three features of Table~\ref{tab:labor-rules} matter for timing. The rules aggregate: the WARN Act counts losses at a site over any 90 days unless their causes are separate and distinct. In Europe consultation precedes the decision, starting once redundancies are contemplated \citep{cjeu2009akavan}, and it covers technology: German works councils are consulted on the use of AI in work processes and must agree to monitoring systems, failing which a conciliation board decides; French councils are consulted on new technologies, for one month or two with an expert; and employers that deploy a high-risk AI system at work must inform workers' representatives first, an obligation of the AI Act that Regulation (EU) 2026/1744 deferred to 2 December 2027 for the employment uses listed in Annex III (2 August 2028 for systems under Annex I) \citep{eu2024regulation,eu2026regulation}. And the rules move: the United Kingdom doubled the maximum protective award in April 2026 and may add an organization-wide threshold \citep{dbt2026threshold}, and India's Industrial Relations Code, in force since 21 November 2025 \citep{pib2025codes}, limits prior permission to factories, mines and plantations, so an offshore software center falls, subject to state law, under the general rules. The same rules price an acquisition that would replace a target's staff: major increases in employment protection cut takeovers by 14 to 27\% and their synergies by over half, more so where deals have more scope to restructure the workforce \citep{dessaint2017employment}.

\begin{proposition}[Lead times and the timing of waves]\label{prop:labor-lead}
Under Assumption~\ref{ass:labor-proc}, take jurisdiction $\mathfrak c$ of a wave (dropped from the notation where clear), its saving $W_{\mathfrak c}$ a year, and a procedure starting at $t^A\ge0$.
\begin{enumerate}
\item[(a)] (Cutover.) Dismissals take effect at $t^{\rm dis}=\max\{t^{\rm gate},\,t^A+T^{\rm leg}\}$. The saving starts $\max\{T^{\rm st},T^{\rm leg}-T_{\rm par}\}$ later than without constraints if the procedure starts at date 0, and $T^{\rm st}+T^{\rm leg}$ later if it starts at the gate.
\item[(b)] (Streams.) Cutting over everywhere at once, at $\bar t:=\max_{\mathfrak c}t^{\rm dis}_{\mathfrak c}$, forgoes $\sum_{\mathfrak c}W_{\mathfrak c}(\bar t-t^{\rm dis}_{\mathfrak c})$. If each separate cutover stream costs the overhead $O_w$ of Assumption~\ref{ass:ext-program}, optimal streams are intervals in the order of $t^{\rm dis}$, and splitting off the fastest group $\mathfrak C_1\subset\mathfrak C$ pays if and only if $\sum_{\mathfrak C_1}W_{\mathfrak c}(\bar t-\max_{\mathfrak C_1}t^{\rm dis}_{\mathfrak c})>O_w$.
\item[(c)] (Windows.) Displacing $N^{\rm dis}$ at one establishment while staying below the threshold in every window takes at least $(\lceil N^{\rm dis}/(N^{\rm thr}-1)\rceil-1)\,T^{\rm win}$.
\item[(d)] (When to start.) Let the announced state cost $c^{\rm ann}$ a year until dismissals or withdrawal, withdrawal after a failure cost $C^{\rm wd}$, and success (probability $\varpi$) be revealed at the gate. The expected cost is minimized at $t^A=(t^{\rm gate}-T^{\rm leg})_+$ or at $t^A=t^{\rm gate}$, and the gate is optimal if and only if
\begin{equation}\label{eq:labor-ann}
\varpi\,WT^{\rm ov}<(1-\varpi)\bigl(c^{\rm ann}T^{\rm ov}+C^{\rm wd}\bigr),\qquad T^{\rm ov}:=\min\{T^{\rm leg},t^{\rm gate}\},
\end{equation}
that is, if $\varpi<\varpi^A:=(c^{\rm ann}T^{\rm ov}+C^{\rm wd})/(WT^{\rm ov}+c^{\rm ann}T^{\rm ov}+C^{\rm wd})$. If the law requires a start by the date $t^{\rm cont}<t^{\rm gate}$ at which redundancies are first contemplated, the optimum is $0$, $t^{\rm cont}$, or $t^{\rm gate}-T^{\rm leg}$ if it lies between them.
\item[(e)] (Free parallel running.) If dismissals cannot take effect until $T^{\rm hold}$ after the run starts, a year of running before $T^{\rm hold}$ costs only $\hat c':=\nu_{\rm rec}E_p+C^{\rm dup}$, since legacy staff are paid and the saving is unavailable until then anyway. The optimal run is
\begin{equation}\label{eq:labor-freerun}
T^{\rm opt}_{\rm par}=\min\bigl\{T^{**}_{\rm par},\ \max\{T^{\rm hold},T^*_{\rm par}\}\bigr\},
\end{equation}
with $T^*_{\rm par}$ from \eqref{eq:tasks-tpar} and $T^{**}_{\rm par}$ from \eqref{eq:tasks-tpar} with $\hat c'$ in place of $\hat c$.
\end{enumerate}
\end{proposition}

\noindent\emph{Intuition and example.} Technology consultation can matter more than the redundancy procedure. For 150 redundancies at five years' tenure the leads are 0.16 years in the United States, 0.12 in the United Kingdom, 0.08 in India, 0.42 in Germany (with three months of consultation, assumed) and 0.47 in France. With a six-month run and assumed technology consultations of three months in Germany and two in France, a procedure started at date 0 gives cutover at 0.5 years in the United States, the United Kingdom and India, 0.67 in France and 0.75 in Germany (0.66, 0.62, 0.58, 1.14 and 1.17 from the gate); a U.S. part saving \$3 million a year forgoes \$0.75 million waiting for Germany. Keeping 200 British redundancies below 20 per 90 days would take 900 days, so the engine counts a program's dismissals at an establishment within a window together. In France, with $W=\$1$ million, $c^{\rm ann}=\$0.3$ million and $C^{\rm wd}=\$0.2$ million, the procedure should start early if success is more likely than 0.42; and with $\hat c=\$600{,}000$, $\hat c'=\$130{,}000$, $(1-\varpi)(C^{\rm str}+L_{\rm roll})=\$550{,}000$ and $T_{\rm det}=0.25$, the French lead extends the optimal run from 0.32 to 0.47 years, cutting the risk of an undetected failure from 6.8\% to 3.8\%.

\subsection{The employees who stay}\label{sec:labor-survivors}

\begin{assumption}[Survivors, leavers and knowledge holders]\label{ass:labor-survivor}
A unit of $N^{\rm hc}$ FTE loses $N^{\rm dis}$ in a wave (the share $\mathsf d:=N^{\rm dis}/N^{\rm hc}$), leaving $N^{\rm sv}=N^{\rm hc}-N^{\rm dis}$ survivors (the residual organization of Definition~\ref{def:ext-org}).
\begin{enumerate}
\item[(i)] From the announcement, survivors quit at the hazard $\nu^q+\nu^s(\mathsf d)$ a year for $T_s$ years instead of $\nu^q$, with $\nu^s(0)=0$ and $\nu^s(x)/x$ nonincreasing; staff due to be displaced quit at $\nu^D\ge\nu^q$ until dismissed.
\item[(ii)] A survivor who quits is replaced at $c^{\rm rep}$ years of the loaded wage $w$ (hiring cost plus the new hire's shortfall), and survivors' output falls by the share $\delta^{\rm eff}(\mathsf d)$ for $T_s$.
\item[(iii)] The wave depends on $n_K$ critical knowledge holders, whose excess hazard $\nu^x$ ($\nu^D-\nu^q$ if they are to be displaced, $\nu^s(\mathsf d)$ otherwise) is a source of departures independent of the baseline in $\varpi^0_p$. A departure during the run causes failure with probability $1-\kappa_{\rm doc}$. A retention award $C^{\rm ret}$, paid at cutover to each holder still employed, removes the share $\kappa_{\rm ret}$ of the excess hazard.
\end{enumerate}
\end{assumption}

The evidence supports the direction of each term more than its size. Downsizing raises voluntary turnover, with an attenuated slope \citep{trevor2008keeping,kim2025preventing}; in 1,620 stores a layoff announcement induced about 2.2 quits per store, against 0.17 after a dismissal \citep{sajjadiani2023who}; and the better workers leave distressed firms first \citep{schwerdt2011labor,lengermann2002abandoning}. Survivors cut performance by 43\% in an experiment \citep{drzensky2016hidden}, sales fell 26\% after closure announcements in a bakery chain \citep{friebel2021downsizing}, and downsizing precedes lower profitability and performance \citep{guthrie2008dumb,park2013turnover}, less so where high-involvement practices are kept \citep{zatzick2006high}. Hiring a skilled worker costs 10 to 17 weeks of wages \citep{blatter2012costs}, and external hires underperform insiders for about two years \citep{bidwell2011paying}. We found no measured rate of internal redeployment, so $u_d$ is a scenario parameter bounded in part (e); retraining effects turn positive after two to three years \citep{card2018what} and are large in screened sectoral programs \citep{katz2022why}.

\begin{proposition}[Survivors, the people plan and success]\label{prop:labor-survivor}
Under Assumption~\ref{ass:labor-survivor}:
\begin{enumerate}
\item[(a)] (Backfill.) The announcement causes $N^{\rm q}=N^{\rm sv}e^{-\nu^qT_s}[1-e^{-\nu^s(\mathsf d)T_s}]\le N^{\rm sv}\nu^s(\mathsf d)T_s$ extra quits, which the people plan must replace, at the cost $C^{\rm surv}=w(c^{\rm rep}N^{\rm q}+N^{\rm sv}\delta^{\rm eff}(\mathsf d)T_s)$, added to $C^{\rm str}$.
\item[(b)] (One cut or several.) To first order, separately announced tranches never lower the extra quits: $\nu^s$ is subadditive, and $n$ equal tranches give $n\nu^s(\mathsf d/n)$, nondecreasing in $n$. If $\nu^s(x)=\eta_0\mathbf 1\{x>0\}+\eta_1x$ and $\delta^{\rm eff}(x)=\delta_0\mathbf 1\{x>0\}+\delta_1x$, each announced wave adds the fixed charge $O_s:=wN^{\rm sv}T_s(c^{\rm rep}\eta_0+\delta_0)$, which enters the wave problem of Proposition~\ref{prop:ext-waves} like the overhead $O_w$.
\item[(c)] (Success.) If the procedure starts $\tau^{\rm ann}$ before the gate, within the run,
\begin{equation}\label{eq:labor-success}
\varpi_p=\varpi^0_p\bigl[\kappa_{\rm doc}+(1-\kappa_{\rm doc})e^{-(1-\kappa_{\rm ret})\nu^x\tau^{\rm ann}}\bigr]^{n_K}\ \ge\ \varpi^0_p\exp\bigl[-n_K(1-\kappa_{\rm doc})(1-\kappa_{\rm ret})\nu^x\tau^{\rm ann}\bigr],
\end{equation}
a structural form of the affected-headcount driver in \eqref{eq:tasks-success}; starting at the gate avoids the loss.
\item[(d)] (Retention awards.) With the values $Y_s$ and $Y_f$ of Proposition~\ref{prop:tasks-pilot}, awards pay if and only if $[\varpi_p(\kappa_{\rm ret})-\varpi_p(0)](Y_s-Y_f)>n_KC^{\rm ret}q_{\rm stay}$, with $q_{\rm stay}=e^{-(\nu^q+(1-\kappa_{\rm ret})\nu^x)\tau^{\rm ann}}$; to first order in $\nu^x\tau^{\rm ann}$, if $C^{\rm ret}<\varpi^0_p(1-\kappa_{\rm doc})\kappa_{\rm ret}\nu^x\tau^{\rm ann}(Y_s-Y_f)$.
\item[(e)] (Redeployment and attrition.) If quits elsewhere open vacancies at the rate $\nu^qN^{\rm oth}$ under a hiring freeze and a displaced worker fills one with probability $\varpi^{\rm rd}$, then over the time $T^w=t^{\rm dis}-t^A$ from announcement to dismissal, $u_d\le\min\{1,\varpi^{\rm rd}\nu^qN^{\rm oth}T^w/N^{\rm dis}\}$, a bound that longer legal leads raise.
\end{enumerate}
\end{proposition}

\noindent\emph{Intuition and example.} The people a run depends on are often those it will displace, and they leave first; the evidence concerns survivors, so $\nu^D$ is a scenario parameter. Reading the 2.2 induced quits per store of about 93.6 workers (as reported in press coverage) as 2.35\% of survivors, a unit with 400 survivors loses 9.4 more people, \$0.53 million at a \$100{,}000 wage and $c^{\rm rep}=0.56$ (13.5 weeks of hiring cost plus an assumed shortfall); as a charge per wave it raises the optimal wave of Proposition~\ref{prop:ext-waves} from 5.29 to 5.67 processes. With three displaced holders, an excess hazard of 0.5, $\kappa_{\rm doc}=0.2$ and the French lead within the run, success falls to 57\% of baseline; awards removing 80\% of the excess hazard restore 89\%, worth \$0.45 million on a \$1.9 million stake at a baseline of 0.75, against \$0.16 million of awards. With quits of 1.1\% a month \citep{bls2026jolts}, 2,000 other staff and 30\% placement, the French lead allows at most 37.5\% of 100 displaced staff to be redeployed.

\subsection{Engineering talent}\label{sec:labor-talent}

\begin{assumption}[Engineering capacity]\label{ass:labor-talent}
Automating a process takes $e_{\rm bld}$ person-years of AI and agent engineering ($l_{\rm AI}F/(w^{\rm AI}t_{\rm FTE})$ in Definition~\ref{def:ext-org}) and running it $e_{\rm mnt}$ FTE a year. Capacity $\Gamma(t)$ is the incumbents $\Gamma_0$ plus hires. A hire starts $T^{\rm lead}=T^{\rm vac}+T^{\rm cn}$ after the requisition (vacancy duration plus the candidate's notice) and delivers $1-\delta^{\rm rp}$ of an FTE in its first $T^{\rm rp}$ years. AI engineers earn $w^{\rm AI}=(1+\mathsf p)w^{\rm SE}$, the loaded software-developer wage plus the premium $\mathsf p$; each hire costs $c^{\rm hire}w^{\rm AI}$. Contractors cost $(1+u_{\rm ct})w^{\rm AI}$ and start after $T^{\rm ct}\le T^{\rm lead}$.
\end{assumption}

U.S. postings requiring AI skills offered 11\% more within the firm and 5\% more within the job title \citep{alekseeva2021demand}, a freelancing platform shows 21\% \citep{stephany2024what}, and AI-exposed establishments cut other hiring \citep{acemoglu2022artificial}; a practitioner analysis reports 28\%, our upper bound \citep{lightcast2025beyond}. The median software developer earned \$135{,}980 in May 2025 \citep{bls2026oewsdev}. Vacancies stay open 29 to 35 days in information and finance \citep{davis2013establishment}, and 35\% of organizations in a self-selected U.K. survey had a hard-to-fill AI vacancy in 2025 \citep{dsit2025ai}. Agent engineering has no such measure, so its premium and lead are scenario parameters.

\begin{proposition}[Talent supply and the pace of the roadmap]\label{prop:labor-talent}
Under Assumption~\ref{ass:labor-talent}:
\begin{enumerate}
\item[(a)] (Premium.) The build costs $F_0(1+l_{\rm AI}\mathsf p)$, with $F_0$ its cost at $w^{\rm SE}$, so the minimum horizon of Proposition~\ref{prop:tasks-transition}(b) rises by $l_{\rm AI}\mathsf pF_0/(\varpi_p\bar b_pW)$.
\item[(b)] (Maintenance ceiling.) With constant capacity $\Gamma$, $e_{\rm bld}\dot N=\Gamma-e_{\rm mnt}N$, so $N(t)=\bar N_\Gamma(1-e^{-e_{\rm mnt}t/e_{\rm bld}})$ with the ceiling $\bar N_\Gamma:=\Gamma/e_{\rm mnt}$, and reaching $N<\bar N_\Gamma$ takes
\begin{equation}\label{eq:labor-TGamma}
T_\Gamma(N)=\frac{e_{\rm bld}}{e_{\rm mnt}}\ln\frac{\bar N_\Gamma}{\bar N_\Gamma-N}=\frac{e_{\rm bld}N}{\Gamma}\cdot\frac{-\ln(1-y)}{y},\qquad y:=\frac N{\bar N_\Gamma}.
\end{equation}
The stretch factor $-\ln(1-y)/y$ rises from 1 to infinity, so capacity caps scope as well as pace.
\item[(c)] (Hiring.) If $\Gamma_0=0$ and the firm hires $\dot\Gamma^{\rm hire}$ FTE a year from date 0, $e_{\rm bld}N$ person-years arrive at $T^{\rm lead}+\sqrt{2e_{\rm bld}N/\dot\Gamma^{\rm hire}}$ without the ramp: the lead adds one for one, the hiring rate through its square root.
\item[(d)] (Hire or contract.) For a gap of $\Delta\Gamma$ FTE lasting $T^{\rm need}\ge T^{\rm rp}$ that holds back savings of $\hat W$ a year, hiring costs $\Delta\Gamma w^{\rm AI}(T^{\rm need}+\delta^{\rm rp}T^{\rm rp}+c^{\rm hire})+\hat WT^{\rm lead}$ and contracting $\Delta\Gamma w^{\rm AI}(1+u_{\rm ct})T^{\rm need}+\hat WT^{\rm ct}$, so contracting is cheaper if and only if
\[
T^{\rm need}<\bigl[\delta^{\rm rp}T^{\rm rp}+c^{\rm hire}+\hat W(T^{\rm lead}-T^{\rm ct})/(\Delta\Gamma w^{\rm AI})\bigr]/u_{\rm ct}.
\]
Build peaks favor contractors; maintenance, a need without end, favors hiring or retraining.
\end{enumerate}
\end{proposition}

\noindent\emph{Example.} With the engine's base wage of \labWSE, the \labPremRate{} premium and the payroll loading of \labTauL{} on an annual wage (June 2026 Employer Costs for Employee Compensation), $w^{\rm AI}=\labWAI$. In the running example (Table~\ref{tab:tasks-example}) the build \icFbuild{} is engineering time priced at the premium, so the premium is \labPremAmt{} of the \icF{} investment, or \labPremShareF{} ($l_{\rm AI}=\labLAI$ and $l_{\rm AI}\mathsf p=\labLAIp$); integration, \icFint, is most of the rest. Part (a) lengthens the minimum horizon only where $W>0$: the example's $W$ is \icW, so no horizon recovers the investment, and the premium only deepens the loss. With $e_{\rm bld}=1.5$ and the engine's $e_{\rm mnt}=0.48$, six engineers sustain at most 12.5 processes and need 10.1 years for 12, not 3.0 (4.16 years, with a ceiling of 24, at $e_{\rm mnt}=0.25$). Hiring ten a year after a quarter-year lead delivers 12 processes at 2.36 years (2.15 without maintenance, the closed form of part (c)). With a 50\% contractor markup, 13.5 weeks of hiring cost, a 15\% shortfall over a hire's first two years, a two-month head start for contractors and \$2 million a year waiting, contract a five-FTE gap lasting under \labContractT{} years.

\subsection{Location and the transfer of task data}\label{sec:labor-location}

\begin{lemma}[Wage index and the adoption date]\label{lem:labor-location}
In the setting of Proposition~\ref{prop:ext-timing}(d)(ii), let a site in $\mathfrak c$ pay the wage index $\mathsf w_{\mathfrak c}\le1$, so that the alternative costs $\mathsf w_{\mathfrak c}c_0$ and the quality cost is $\mathsf w_{\mathfrak c}c^Q_{\rm lab}+c^Q_{\rm loss}$, while AI spend and the build do not scale. With $\bar c:=c_0-c^Q_{\rm lab}$ and $c^{\rm fx}:=c^Q_{\rm loss}+rF_0/v\ge0$, if both adoption dates are interior the site adopts
\[
T^*_{\mathfrak c}-T^*_1=\frac1{\delta_\pi}\ln\frac{\bar c-c^{\rm fx}}{\mathsf w_{\mathfrak c}\bar c-c^{\rm fx}}\ \ge\ \frac{\ln(1/\mathsf w_{\mathfrak c})}{\delta_\pi}
\]
years later, with equality if and only if $c^{\rm fx}=0$; if $\mathsf w_{\mathfrak c}\bar c\le c^{\rm fx}$, automation never pays there.
\end{lemma}

At market exchange rates and relative to the 2024 U.S. average wage, the index is 0.93 for Germany, 0.66 for France, 0.33 for Poland and 0.042 for India \citep{oecd2025pensions}, so under the lemma, at a fivefold annual price decline, an Indian site adopts at least 1.97 years after a U.S. one. The engine prefers roster wages. O*NET time shares transfer imperfectly: across 47 countries job content differs substantially within occupations, and occupational structure explains about a fifth of the variation in routine task intensity \citep{lewandowski2022technology}; returns to tasks vary with institutions \citep{delarica2020job}; and O*NET and its crosswalks have known gaps \citep{handel2016onet,bls2012isco,pizzinelli2023labor}. The engine therefore widens time-share uncertainty outside the United States.

\subsection{Data residency}\label{sec:labor-residency}

\begin{assumption}[Residency and transfer]\label{ass:labor-residency}
The AI system for activity $j$ runs in a region $\mathfrak g$ (in the country, in the region, or abroad). The rules of the data subjects' jurisdiction make the share $\omega^{\rm res}_j(\mathfrak g)\in[0,1]$ of the inputs unusable at $\mathfrak g$, allow a further share only through a transfer mechanism, and leave the rest free. A mechanism costs $F^{\rm tr}_{\mathfrak g}$ a year and lapses at the hazard $\nu^{\rm inv}_{\mathfrak g}$, after which replacing it or moving the workload costs $C^{\rm sw}$. Region $\mathfrak g$ multiplies AI spend by $1+p^{\rm reg}_{\mathfrak g}$ and lags frontier capability by $T^{\rm lag}_{\mathfrak g}$.
\end{assumption}

The EU's General Data Protection Regulation allows transfers to countries found adequate, otherwise under safeguards such as standard contractual clauses, and otherwise only under narrow derogations \citep{eu2016regulation,eu2021scc}; China and India hold no adequacy decision \citep{ec2026adequacy}. U.S. transfers rely on the Data Privacy Framework of July 2023 \citep{eu2023dpf}, upheld at first instance in September 2025 and under appeal \citep{eugc2025latombe}, whose two predecessors were invalidated \citep{cjeu2015schrems,cjeu2020schrems}; a June 2026 Supreme Court holding on the removal of Federal Trade Commission members \citep{us2026slaughter} raises a question about the independence of the oversight it relies on. China requires an assessment, certification or standard contract for exports of personal information, with assessment above one million people (10,000 for sensitive data) a year, and local storage for critical infrastructure \citep{cn2021pipl,cn2024flows,cn2025certification}. India may restrict transfers to notified countries, under rules mostly in force from May 2027 \citep{in2023dpdp,in2025dpdprules}.

\begin{proposition}[Residency limits on readiness and scope]\label{prop:labor-residency}
Under Assumption~\ref{ass:labor-residency}:
\begin{enumerate}
\item[(a)] (Readiness cap and permissibility.) At region $\mathfrak g$ readiness is at most $\bar\iota_j(\mathfrak g):=(1-\omega^{\rm res}_j(\mathfrak g))\iota_j$. The AI modes that use the restricted inputs are feasible at $\mathfrak g$ if and only if $\bar\iota_j(\mathfrak g)\ge\underline\iota_j$ and the capability test of Assumption~\ref{ass:ext-capability} holds at $t-T^{\rm lag}_{\mathfrak g}$, and permissible if and only if some region passes both (else $j$ leaves $\mathcal J^{\rm per}$).
\item[(b)] (Process threshold.) Under the conditions of Proposition~\ref{prop:tasks-readiness}, with $\iota_p\ge\iota^*_p$ and restrictions acting in proportion on every activity, full automation at $\mathfrak g$ stays optimal if and only if $\omega^{\rm res}_p(\mathfrak g)\le\bar\omega^{\rm res}_p:=1-\iota^*_p/\iota_p$; above it every activity reverts to its best non-automated mode. The value $[W_p(\mathcal J_p;(1-\omega)\iota_p)]_+$ falls continuously to zero at $\omega=\bar\omega^{\rm res}_p$ if $\iota^*_p>\underline\iota_p$, and drops by $W_p(\mathcal J_p;\underline\iota_p)>0$ there if the feasibility floor binds. One activity pushed below its floor makes full automation infeasible.
\item[(c)] (Choice of region.) The best region minimizes
\begin{equation}\label{eq:labor-region}
c_{\mathrm A,j}(\mathfrak g)=(1+p^{\rm reg}_{\mathfrak g})\pi n_j+\hat c^Q_j(\bar\iota_j(\mathfrak g))+\bigl(F^{\rm tr}_{\mathfrak g}+\nu^{\rm inv}_{\mathfrak g}C^{\rm sw}\bigr)/v_j
\end{equation}
plus costs common to all regions. The premium $\Delta c^{\rm res}_j:=\min_{\mathfrak g}c_{\mathrm A,j}(\mathfrak g)-c_{\mathrm A,j}(\text{no restriction})\ge0$ lowers the automation advantage $\mathcal B_j$ of Definition~\ref{def:ext-effcost} by up to $\Delta c^{\rm res}_j$, and the activity leaves the economical universe through mode A if $0<\mathcal B_j\le\Delta c^{\rm res}_j$ and no other AI mode wins.
\end{enumerate}
\end{proposition}

\noindent\emph{Example.} A process at readiness 0.9 with a threshold of 0.85 tolerates restrictions on only 5.6\% of its inputs; restricting 10\% returns the whole chain to augmentation. Two EU--U.S. frameworks lapsed in about 22.5 framework-years, a crude hazard of 0.089 a year (\$44{,}000 a year on a \$0.5 million switching cost), a scenario range in the engine.

\subsection{Pace}\label{sec:labor-pace}

\begin{corollary}[Pace under legal, people and talent constraints]\label{cor:labor-pace}
Let a program of $N$ processes run in sequential waves, each lasting $T^{\rm eff}:=\max_{\mathfrak c}t^{\rm dis}_{\mathfrak c}$ (Proposition~\ref{prop:labor-lead}(a), with the start chosen by part (d)). Propositions~\ref{prop:ext-waves} and \ref{prop:ext-bandwidth} hold with $T_{\rm par}$ replaced by $T^{\rm eff}$ and $O_w$ by $O_w+O_s$, and with engineering capacity at most $\Gamma_{\max}$,
\[
T_{\rm prog}\ \ge\ \max\Bigl\{\frac{Na}{K_{\rm chg}},\ \frac{NeT^{\rm eff}}{\hat\varkappa\Phi},\ T_{\Gamma_{\max}}(N)\Bigr\},
\]
which is finite only if $N<\Gamma_{\max}/e_{\rm mnt}$. In the example of Proposition~\ref{prop:ext-waves}, six engineers at $e_{\rm mnt}=0.48$ need 10.1 years, close to their ceiling of 12.5 processes, and bind long before change capacity (1.5 years) and bandwidth (1.28); at $e_{\rm mnt}=0.25$ they still bind, at 4.16 years.
\end{corollary}

